\section{Билет 11. Матрица проектирования. Метод Розена}

\ticketq{Матрица проектирования. Метод проекции градиента Розена.}

\paragraph{Постановка и идея.}
\[
  f(x) \to \min , \qquad Ax = b , \qquad Cx \le d ,
\]
$f$ дифференцируема. В безусловной задаче направление убывания —
$-\nabla f(x_k)$, но оно может вывести за пределы допустимого множества.
Поэтому антиградиент \emph{проектируют} на множество допустимых направлений:
из $A(x + \alpha s) = b$ и $Ax = b$ следует $As = 0$, то есть допустимые
направления лежат в подпространстве $\{s \mid As = 0\}$.

\paragraph{Матрица проектирования.}
$P \in \R^{n\times n}$ — \textbf{матрица проектирования}, если
$P\T = P$ и $P^2 = P$.

\textbf{Лемма (общие свойства).} \textit{$Q = E - P$ — тоже проектор;
$P, Q \succeq 0$; $\operatorname{Im}P \perp \operatorname{Im}Q$; разложение
$x = Px + Qx$ единственно.}

\textsc{Док-во.} $Q\T = E - P\T = Q$ и
$Q^2 = E - 2P + P^2 = E - P = Q$. Далее
$x\T Px = x\T P\T Px = \norm{Px}^2 \ge 0$. Ортогональность:
$(Px_1)\T Qx_2 = x_1\T(P - P^2)x_2 = 0$. Единственность: если
$x = y+z = y_1+z_1$ с $y,y_1 \in L = \operatorname{Im}P$,
$z,z_1 \in L' = \operatorname{Im}Q$, то
$y - y_1 = z_1 - z \in L \cap L' = \{0\}$. $\blacksquare$

Для матрицы активных ограничений $A \in \R^{m\times n}$ с $\rank A = m$
(тогда $AA\T$ невырождена) проектор на $\{s \mid As = 0\}$:
\[
  P = E - A\T\big(AA\T\big)^{-1}A .
\]

\paragraph{Теорема (свойства проектора).}
\textit{1) $P\T = P$; 2) $P^2 = P$; 3) $AP = 0$; 4) $Az = 0 \Rightarrow Pz = z$.}

\textsc{Док-во.}
1) $P\T = E\T - A\T\big((AA\T)^{-1}\big)\T A$; $AA\T$ симметрична
$\Rightarrow$ её обратная симметрична $\Rightarrow$ $P\T = P$.

2) $P^2 = E - 2A\T(AA\T)^{-1}A + A\T(AA\T)^{-1}\underbrace{AA\T(AA\T)^{-1}}_{=\,E}A
   = E - A\T(AA\T)^{-1}A = P$.

3) $AP = A - \underbrace{AA\T(AA\T)^{-1}}_{=\,E}A = A - A = 0$.

4) $Pz = z - A\T(AA\T)^{-1}\underbrace{Az}_{=\,0} = z$. $\blacksquare$

\paragraph{Направление метода.}
$g = \nabla f(x)$, $s = -Pg$.

\textbf{Теорема (направление спуска).} \textit{$Pg \ne 0 \Rightarrow g\T s < 0$.}

\textsc{Док-во.} $P\T = P$, $P^2 = P$ $\Rightarrow$
\[
  g\T s = -g\T P g = -g\T P\T P g = -(Pg)\T(Pg) = -\norm{Pg}^2 < 0 .
  \ \blacksquare
\]

\textbf{Теорема (возможное направление).} \textit{$Ax = b$, $s = -P\nabla f(x)$
$\Rightarrow$ $A(x + \alpha s) = b$.}

\textsc{Док-во.} $As = -AP\nabla f(x) = 0$ по свойству 3 $\Rightarrow$
$A(x+\alpha s) = Ax + \alpha\cdot 0 = b$. $\blacksquare$

\paragraph{Условие оптимальности и множители Лагранжа.}
Пусть $Pg = 0$. Тогда
\[
  g - A\T(AA\T)^{-1}Ag = 0
  \ \Rightarrow\ g = A\T\mu , \quad \mu = (AA\T)^{-1}Ag ,
\]
то есть градиент цели выражается через строки активных ограничений — это
условие Лагранжа $\nabla f(x) + A\T\lambda = 0$ при $\lambda = -\mu$. Для
ограничений вида $a_i\T x \le b_i$ условия Куна--Таккера требуют
$\lambda_i \ge 0$:
\begin{itemize}[nosep,left=0pt]
  \item все $\lambda_i \ge 0$ — точка удовлетворяет необходимым условиям
        оптимальности, конец;
  \item $\exists\,\lambda_i < 0$ — это ограничение исключают из активного
        набора и строят проектор заново.
\end{itemize}

\textbf{Почему исключение работает.} Пусть $u_j < 0$ (множитель при активном
неравенстве $f_j\T x \le g_j$). Уберём строку $f_j$ из $A$, получив $A_1$ и
$P_1$. Тогда $\nabla f = -A_1\T\omega_1 - u_j f_j\T$ и
\[
  P_1\nabla f = -u_j P_1 f_j\T \ne 0
\]
(иначе $u_j = 0$ из линейной независимости строк). Новое направление
$s' = -P_1\nabla f = u_j P_1 f_j\T$ одновременно:
\[
  \nabla\T f\,s' = -\norm{P_1\nabla f}^2 < 0 , \qquad
  A_1 s' = 0 , \qquad
  f_j s' = u_j\norm{P_1 f_j\T}^2 < 0 ,
\]
то есть это направление спуска, сохраняющее остальные активные ограничения и
уводящее \emph{внутрь} по снятому неравенству — оно перестаёт быть активным, и
движение допустимо.

\paragraph{Выбор шага.}
Активный набор в $x_k$: $I_k = \{i \mid a_i\T x_k = b_i\}$, $A_k$ — матрица
его строк, $P_k = E - A_k\T(A_kA_k\T)^{-1}A_k$, $s_k = -P_k g_k$. Для
неактивных ограничений нужно $a_i\T(x_k + \alpha s_k) \le b_i$, откуда
\[
  \alpha_{\max} = \min_{i \notin I_k,\ a_i\T s_k > 0}
    \frac{b_i - a_i\T x_k}{a_i\T s_k}
  \qquad (\alpha_{\max} = +\infty, \text{ если таких } i \text{ нет}) ,
\]
и $\alpha_k = \argmin_{0 \le \alpha \le \alpha_{\max}} f(x_k + \alpha s_k)$.

\paragraph{Алгоритм.}
\begin{csteps}
  \item Задать допустимую $x_0$, $\eps > 0$.
  \item Определить активный набор $I_k$, построить $A_k$ и $P_k$ (если
        активных нет — $P_k = E$).
  \item $g_k = \nabla f(x_k)$, $s_k = -P_k g_k$.
  \item $\norm{s_k} > \eps$ — направление найдено, к~шагу~5. Иначе вычислить
        множители Лагранжа: все $\lambda_i \ge 0$ — конец; иначе удалить
        ограничение с $\lambda_i < 0$ из $I_k$ и к~шагу~2.
  \item Найти $\alpha_{\max}$ и $\alpha_k$; $x_{k+1} = x_k + \alpha_k s_k$.
  \item Если шаг сделал новое ограничение активным — добавить его в $I_k$;
        $k := k+1$, к~шагу~2.
\end{csteps}

\paragraph{Ограничения.}
Нужна допустимая начальная точка; ограничения должны быть линейными (или
локально линеаризуемыми); при линейно зависимых строках активных ограничений
$A_kA_k\T$ вырождена; при $P_kg_k = 0$ нужна дополнительная проверка
множителей; для невыпуклой $f$ — локальный минимум.

\seefull{гл.~3, \S5--\S6 (численные примеры там же)}
